\noindent\hspace*{2em}神经网络计算图的多核执行能够提高计算并行度，而任务切分带来的跨域搬运、片上容量占用与共享带宽竞争会影响总体工期。

本文研究独立任务、同核复用与共享只读缓存三种场景下的多核切图与调度。首先通过\textbf{计算依赖提取与弱连通分量分析}，构造子图划分与核心分配方案；接着根据张量的消费范围和物理副本生存区间，建立结构搬运与容量约束模型；最后引入 \textbf{FIFO 缓存事件模型}，由题给 Python 事件模拟程序评价 100 张计算图在 1 至 5 核下的调度结果。

\textbf{针对问题一}，建立 \textbf{Task 驻留域结构搬运模型}。首先提取计算依赖及张量的生产、消费关系，按接收 Task 统计输入读入、跨任务搬运与最终写回，并根据同核次序和前驱依赖确定 Task 激活条件。接着以\textbf{整图分配、分量装箱与拓扑切块}生成候选，用静态评分安排评价顺序，再由事件评估程序按 Makespan 优先、额外搬运次之选定方案；对每张图在不超过目标核数的已评价计划及整图基准中取最短工期。相对整图单核基准，五核平均加速比为 \textbf{3.4132}；图 001 的工期由 233110 周期缩短至 \textbf{47502 周期}，新增搬运量为 \textbf{4608 字节}。

\textbf{针对问题二}，建立\textbf{核心驻留与物理副本容量模型}。首先将同核子图合并为核级 Task，按接收核心统计跨核搬运。其次根据物理副本从空间申请至最后读者完成的存活区间，检查 L1、UB 的同时占用，分别核算结构搬运与换出恢复。然后以\textbf{关键路径优先}和释放优先生成拓扑切块候选，结合贪心分核与局部调整，经题给事件评价选定方案。采用逐图继承规则后，五核平均加速比为 \textbf{3.3862}；由四核升至五核时，\textbf{65 图}工期缩短、\textbf{35 图}持平，且没有图慢于整图单核基准。

\textbf{针对问题三}，建立\textbf{基于查询与完成事件的 FIFO Cache 模型}。首先按照发射时查询、读取完成后插入和先进先出淘汰的规则，把命中读与未命中读分别纳入 Cache 和 DDR 的独立带宽服务。接着以问题二最终计划为起点，通过共享输入聚合、迁块和核内调序构造候选。然后固定计划比较有无 Cache 的工期，再在相同 Cache 配置下比较调整前后的计划，区分硬件作用与调度适配。采用同一逐图继承规则后，五核相对整图单核基准的平均加速比为 \textbf{3.4272}；\textbf{同图同核数}的 B/C 综合工期比均值为 \textbf{1.0136}。

固定问题二的五核计划比较有无 Cache，100 张图中 \textbf{54 图}工期缩短、\textbf{46 图}持平；在相同 Cache 配置下调整计划，\textbf{7 图}进一步改善。六张代表图的 Cache 容量与带宽扰动中，重新比较已评候选后均保留原计划。

\noindent{\zihao{3}\textbf{关键词：}}多核调度；张量驻留域；拓扑切块；容量约束；FIFO 缓存
